\section{A circular version of $\muLJ$}
\label{sec:cmuLJ-prelims}

In this section we shall develop a variation of $\muLJ$ that does not have rules for (co)iteration, but rather devolves such work to the proof structure.
First, let us set up the basic system of rules we will work with:
\begin{definition}
[$\muLJ$ `without (co)iteration']
Write $\muLJnorec$ for the system of all rules in Figures~\ref{fig:sequent-calculus-mulj}, \ref{fig:unfolding-rules} and \ref{fig:unfolding-further-rules} (but not \ref{fig:co-iteration-rules}).
\end{definition}

% \begin{remark}
% [$\muLJnorec$ as a subsystem of $\muLJ$]
% Let us point out that, at the level of derivability, $\muLJnorec$ is certainly a subsystem of $\muLJ$, in light of Example~\ref{ex:post-fixed-in-muLJ}.
% On the other hand the relationship between them computationally is a little more subtle, as we shall see.
% \anupam{do we need this remark, esp here?}
% \end{remark}

\subsection{`Non-wellfounded' proofs over $\muLJnorec$}

`Coderivations' are generated \emph{coinductively} by the rules of a system, dually to derivations that are generated inductively.
I.e.\ they are possibly infinite proof trees generated by the rules of a system.
% We say that a coderivation is \emph{regular} (or \emph{circular}) if it has only finitely many distinct sub-coderivations.


\begin{defn}
[Coderivations]
A ($\muLJnorec$-)\emph{coderivation} $\der$ is a possibly infinite {rooted} tree (of height $\leq \omega$) generated by the rules of $\muLJnorec$.
Formally, we identify $\der$ with a prefix-closed subset of $ \{0,1\}^*$ (i.e.\ a binary tree) where each node  is labelled by an inference step from $\muLJnorec$ such that, whenever $\alpha\in \{0,1\}^*$ is labelled by a step $\vliiinf{}{}{S}{S_1}{\cdots}{S_n}$, for $n\leq 2$, $\alpha$ has $n$ children in $\der$ labelled by steps with conclusions $S_1, \dots, S_n$ respectively. 

We say that a coderivation is \emph{regular} (or \emph{circular}) if it has only finitely many distinct sub-coderivations.
\end{defn}


A regular coderivation can be represented as a finite labelled graph (possibly with cycles) in the natural way.



\begin{figure}
     \centering
 \[
 \small
   % \vlinf{\rr \mu}{}{\Gamma \seqar \mu X. \sigma}{\Gamma \seqar \sigma(\mu X. \sigma)}
   % \qquad 
   \vlinf{\mulunf}{}{\mygreen{\Gamma}, \red{\mu X. \sigma} \seqar \cyan{\tau}}{\mygreen{\Gamma}, \red{\sigma(\mu X.\sigma)} \seqar \cyan{\tau}}
   \qquad 
   \vlinf{\nurunf}{}{\mygreen{\Gamma} \seqar \red{\nu X.\sigma}}{\mygreen{\Gamma} \seqar \red{\sigma(\nu X.\sigma)}}
   % \qquad
   % \vlinf{\lr \nu}{}{\Gamma, \nu X.\sigma \seqar \tau}{\Gamma, \sigma(\nu X. \sigma)\seqar \tau}
   \]
    \caption{Further unfolding rules for $\mu$ and $\nu$.}
    \label{fig:unfolding-further-rules}
\end{figure}



\subsection{Computing with coderivations}

Just like for usual derivations, the underlying notion of computation for coderivations is cut-reduction, and the notion of representability remains the same.
However we must also adapt the theory of cut-reduction to the different fixed point rules of $\muLJnorec$.

\begin{definition}
    [Cut-reduction on coderivations]\label{defn:rrmu-vs-llmu-cmull}
    $\cutrednorec$ is the smallest relation on $\muLJnorec$-coderivations including all the usual cut-reductions of $\LJ$ and the cut-reductions in Figure~\ref{fig:cut-reduction-cmulj-least-greatest},\footnote{Again, we allow these reductions to be applied on sub-coderivations.}.
% {(see Appendix~\ref{app:simple-types-with-fixed-points}, point~\ref{enum:app2})}.
% 
When speaking of (subsets of) coderivations as computational models, we typically mean with respect to $\cutrednorec^*$.
\end{definition}

\begin{figure}
    \centering
\[
    \footnotesize
           \vlderivation{
\vliin{\cut}{}{\Gamma, \Delta \seqar \tau}{
    \vlin{\rr \mu}{}{\Gamma \seqar \mu X \sigma(X) }{\vlhy{\Gamma \seqar \sigma (\mu X \sigma(X))}}
}{
    \vlin{\mulunf }{}{\Delta , \mu X \sigma (X)\seqar \tau}{\vlhy{\Delta, \sigma(\mu X \sigma(X)) \seqar \tau}}
}
}
  \quad \rewrite \quad      
  \vliinf{\cut}{}{\Gamma, \Delta \seqar \tau}{\Gamma \seqar \sigma(\mu X \sigma(X))}{\Delta, \sigma(\mu X \sigma(X)) \seqar \tau}
  \]
  \[
  \footnotesize
   \vlderivation{
\vliin{\cut}{}{\Gamma, \Delta \seqar \tau}{
    \vlin{\nurunf }{}{\Gamma \seqar \nu X \sigma(X) }{\vlhy{\Gamma \seqar \sigma (\nu X \sigma(X))}}
}{
    \vlin{\lr \nu }{}{\Delta , \nu X \sigma (X)\seqar \tau}{\vlhy{\Delta, \sigma(\nu X \sigma(X))}}
}
}
 \quad \rewrite \quad 
 \vliinf{\cut}{}{\Gamma, \Delta \seqar \tau}{\Gamma \seqar \sigma(\nu X \sigma(X))}{\Delta, \sigma(\nu X \sigma(X)) \seqar \tau}
\]    \caption{Cut-reduction for least and greatest fixed points in $\cmuLJ$.}
    \label{fig:cut-reduction-cmulj-least-greatest}
\end{figure}


% It is implicit in all circumstances in this work that subclasses of coderivations we consider are closed under $\cutrednorec$.
% In particular, let us point out that the class of regular coderivations is closed under $\cutrednorec$. \anupam{say more?}

\begin{example}
    [Decomposing the (co)iterators]
    \label{ex:(co)iter-to-cycles}
    The `(co)iterator' rules of Figure~\ref{fig:co-iteration-rules} can be expressed by regular coderivations using only the unfolding rules for fixed points as follows:
    \begin{equation}
    \label{eq:mul-as-circ-coder}
    %     \vlderivation{
    % \vliin{\cut}{}{\Gamma, \Delta, \mu X \sigma(X) \seqar \tau}{
    % \vlin{\mulunf}{\bullet}{\Gamma, \mu X \sigma(X) \seqar \rho}{
    % \vliin{\cut}{}{\Gamma, \sigma(\mu X \sigma(X))\seqar \rho}{
    % \vliq{\sigma}{}{\Gamma, \sigma(\mu X \sigma(X)) \seqar \sigma(\rho)}{
    % \vlin{\mulunf}{\bullet}{\Gamma, \mu X \sigma(X) \seqar \rho}{\vlhy{\vdots}}
    % }
    % }{
    % \vlhy{\Gamma, \sigma(\rho)\seqar \rho}
    % }
    % }
    % }{
    % \vlhy{\Delta, \rho \seqar \tau}
    % }
    % }
    \small
       \vlderivation{
    \vliin{\cut}{}{\Gamma, \Delta, \mu X \sigma(X) \seqar \tau}{
    \vlin{\mulunf}{\bullet}{\Gamma, \mu X \sigma(X) \seqar \rho}{
    \vlin{\contr}{}{\Gamma,  \sigma(\mu X \sigma(X))\seqar \rho}
    {
    \vliin{\cut}{}{\Gamma,\Gamma,  \sigma(\mu X \sigma(X))\seqar \rho}{
    \vliq{\sigma}{}{\Gamma, \sigma(\mu X \sigma(X)) \seqar \sigma(\rho)}{
    \vlin{\mulunf}{\bullet}{\Gamma, \mu X \sigma(X) \seqar \rho}{\vlhy{\vdots}}
    }
    }{
    \vlhy{\Gamma, \sigma(\rho)\seqar \rho}
    }
    }
    }
    }{
    \vlhy{\Delta, \rho \seqar \tau}
    }
    }
    \end{equation}
    Here we mark with $\bullet$ roots of identical coderivations, a convention that we shall continue to use throughout this work.
    Dually for the coiterator:
    \[
\small
       \vlderivation{
    \vliin{\cut}{}{\Gamma, \Delta \seqar \nu X \sigma(X)}
    {
    \vlhy{\Delta \seqar \tau}
    }
    {
    \vlin{\nurunf}{\bullet}{\Gamma, \tau \seqar \nu X \sigma(X)}{
    \vliq{\contr}{}{\Gamma,  \tau\seqar \sigma(\nu X \sigma(X))}
    {
    \vliin{\cut}{}{\Gamma,\Gamma,  \tau\seqar \sigma(\nu X \sigma(X)}
    {
    \vlhy{\Gamma, \tau\seqar \sigma(\tau)}
    }
    {
    \vliq{\sigma}{}{\Gamma, \sigma(\tau) \seqar \sigma(\nu X \sigma(X))}{
    \vlin{\nurunf}{\bullet}{\Gamma, \tau \seqar \nu X \sigma(X)}{\vlhy{\vdots}}
    }
    }
    }
    }
    }
    }
\]

    Moreover, one can verify that this embedding gives rise to a bona fide simulation of $\cutred^*$ by $\cutrednorec^*$. 
    We do not cover the details at this point, but make a stronger statement later in Proposition~\ref{prop:cmulj-simulates-mulj}.
\end{example}


\begin{example}
    [Functors and $\eta$-expansion of identity]
    Thanks to the decomposition of (co)iterators above, we can derive `functors' in $\cmuLJ$, cf.~Definition~\ref{def:functors-in-muLJ}.
    This gives rise to an `$\eta$-expansion' of identity steps, reducing them to atomic form. 
The critical fixed point cases are:
    \[
    \small
    \vlderivation{
    \vlin{\mulunf}{\bullet}{\mu X \sigma(X)\seqar \mu X \sigma(X)}{
    \vlin{\rr \mu}{}{\sigma(\mu X(\sigma(X))\seqar \mu X \sigma(X)}{
    \vliq{\sigma}{}{\sigma(\mu X \sigma(X))\seqar \sigma(\mu X \sigma(X))}{
    \vlin{\lr \mu}{\bullet}{\mu X \sigma(X)\seqar \mu X \sigma(X)}{\vlhy{\vdots}}
    }
    }
    }
    }
    \qquad
    \small
    \vlderivation{
    \vlin{\lr \nu}{\bullet}{\nu X \sigma(X)\seqar \nu X \sigma(X)}{
    \vlin{\rr {\nu'}}{}{\sigma(\nu X(\sigma(X))\seqar \nu X \sigma(X)}{
    \vliq{\sigma}{}{\sigma(\nu X \sigma(X))\seqar \sigma(\nu X \sigma(X))}{
    \vlin{\lr \nu}{\bullet}{\nu X \sigma(X)\seqar \nu X \sigma(X)}{\vlhy{\vdots}}
    }
    }
    }
    }
    \]
    % 
Notice that the functors $\sigma$ indicated above will depend on smaller identities, cf.~\Cref{def:functors-in-muLJ}, calling the inductive hypothesis.
    % 
    Note that the coderivations above are `logic-independent', and indeed this reduction is common in other circular systems for fixed point logics, such as the modal $\mu$-calculus and $\muMALL$ (see, e.g.,~\cite{BaeldeDS16Infinitaryprooftheory}).
\end{example}



\subsection{A totality criterion}
We shall adapt to our setting a well-known `termination criterion' from non-wellfounded proof theory. First, let us recall some standard  proof theoretic concepts about (co)derivations, similar to those in \cite{BaeldeDS16Infinitaryprooftheory,KuperbergPP21systemT,Das21:CT-preprint,Das21:CT-fscd}.

\begin{defn}
[Ancestry]
\label{defn:ancestry}
Fix a $\muLJnorec$-coderivation $\der$. We say that a type occurrence $\sigma$ is an \emph{immediate ancestor} of a type occurrence $\tau$ in $\der$ if they are types in a premiss and conclusion (respectively) of an inference step and, as typeset in~\Cref{fig:sequent-calculus-mulj},~\Cref{fig:unfolding-rules} and~\Cref{fig:unfolding-further-rules}, have the same colour.
If $\sigma$ and $\tau$ are in some $\mygreen \Gamma$ or $\orange{ \Delta}$, then furthermore they must be in the same position in the list.
\end{defn}




Being a binary relation, immediate ancestry forms a directed graph upon which our correctness criterion is built. 
Our criterion is essentially the same as that from \cite{BaeldeDS16Infinitaryprooftheory}, only for $\muLJ$ instead of $\muMALL$. 



\begin{defn}
    [Threads and progress]\label{defn:progressiveness}
    % \todo{define immediate ancestry}
    A \emph{thread} along (a branch of) $\der$ is a maximal path in $\der$'s graph of immediate ancestry.
We say a thread is \emph{progressing} if it is infinitely often principal and has a smallest infinitely often principal formula that is either a $\mu$-formula on the LHS or a $\nu$-formula on the RHS.
% 
A coderivation $\der$ is \emph{progressing} if each of its infinite branches has a progressing thread.
\end{defn}

We shall use several properties of (progressing) threads in Section~\ref{sec:totality} 
which are relatively standard, e.g.~\cite{Koz83:results-on-mu,Studer08:mu-calc,KupMarVen22:graph-reps-mu-forms}.

\begin{defn}
    [Circular system]
    $\cmuLJ$ is the class of regular progressing $\muLJnorec$-coderivations.
\end{defn}


Referencing Example~\ref{ex:(co)iter-to-cycles}, and for later use, we shall appeal to the notion of \emph{simulation} for comparing models of computation in this work. 
Recalling that we construe $\muLJ$ as a model of computation under $\cutred^*$ and $\cmuLJ$ as a model of computation under $\cutrednorec^*$, we have:

\begin{proposition}[Simulation]
\label{prop:cmulj-simulates-mulj}
$\cmuLJ$ simulates $\muLJ$.
\end{proposition}

\begin{proof}[Proof sketch]
    Replace each instance of a (co)iterator by the corresponding regular coderivation in Example~\ref{ex:(co)iter-to-cycles}.
    Note that those coderivations are indeed progressing due to the progressing thread on $\mu X \sigma(X)$ along the unique infinite branch in the case of $\lr \mu$ (dually for $\rr \nu$).
    The statement follows by closure of $\cmuLJ$ under its rules.
\end{proof}


% \anupam{make a remark about whether $\cutrednorec^*$ is simulated by $\cutred^*$ for derivations, and this reduces to a form of extensionality. We have distinguished them by design choice, but could have chosen to eta-expand identities via explicit functor rules to collapse them. Point out that this will (unsurprisingly) not affect representability, cf.~Luckhardt.}


\begin{figure}
    \centering
      \[\small
  \vlinf{\Nzero}{}{\seqar \N}{}
  \qquad 
  \vlinf{\Nsucc}{}{\mygreen \Gamma \seqar \cyan \N}{\mygreen \Gamma \seqar \cyan \N} 
  \qquad 
  \vliinf{\Ncnd}{}{\mygreen{\Gamma}, \red{\N} \seqar \cyan \sigma}{\mygreen\Gamma \seqar \cyan \sigma}{\mygreen \Gamma, \red \N \seqar \cyan \sigma}
    \]
    \[
    \small
      \vlderivation{
      \vliin{\cut}{}{\Gamma \seqar \sigma}
       {
       \vlin{\rr\N^0}{}{\seqar \N}{\vlhy{}}
       }
      {
 \vliin{\Ncnd}{}{\Gamma, \N \seqar \sigma}{
    \vltr{\derr}{\Gamma \seqar \sigma}{\vlhy{\ }}{\vlhy{}}{\vlhy{\ }}
    }{
    \vltr{\derrr}{\Gamma, \N \seqar \sigma}{\vlhy{\ }}{\vlhy{}}{\vlhy{\ }}
    }
      }
      }
      \quad \cutrednorec \  
      \vlderivation{\vltr{\derr}{\Gamma \seqar \sigma}{\vlhy{\ }}{\vlhy{}}{\vlhy{\ }}} 
      \]
      \[
          \hspace{-1em}
      \small
      \vlderivation{
      \vliin{\cut}{}{\Gamma \seqar \sigma}
       {
       \vlin{\Nsucc}{}{\Gamma \seqar \N}{\vltr{\der}{\Gamma \seqar \N}{\vlhy{\ }}{\vlhy{}}{\vlhy{\ }} }
       }
      {
 \vliin{\Ncnd}{}{\Gamma, \N \seqar \sigma}{\vltr{\derr}{\Gamma \seqar \sigma}{\vlhy{\ }}{\vlhy{}}{\vlhy{\ }}}{\vltr{\derrr}{\Gamma, \N \seqar \sigma}{\vlhy{\ }}{\vlhy{}}{\vlhy{\ }}}
      }
      }   
      \ \cutrednorec
      \vlderivation{
      \vliin{\cut}{}{\Gamma \seqar \sigma}
        {
        \vltr{\der}{\Gamma \seqar \N}{\vlhy{\ }}{\vlhy{}}{\vlhy{\ }}
        }{
      \vltr{\derrr}{\Gamma, \N \seqar \sigma}{\vlhy{\ }}{\vlhy{}}{\vlhy{\ }}        
        }
      }
    \]
    \caption{Native inference rules and cut-reduction steps for $\N$ in $\muLJnorec$.}
    \label{fig:native-n-rules+cutreds-muljnorec}
\end{figure}


\begin{example}
    [Revisiting natural number computation]
    Just like for $\muLJ$, we give native rules for $\N$ in $\muLJnorec$, along with corresponding cut-reductions in Figure~\ref{fig:native-n-rules+cutreds-muljnorec}. We just show how to derive the conditional:
 \[
% \Ncnd\ \dfn \ 
\vlderivation
{
\vlin{\lr {\mu'}}{}{\Gamma, \N \seqar \sigma}
{
\vliin{\lr{+}}{}{\Gamma, \unit + \N \seqar \sigma}
{
\vlin{\lr \unit}{}{\Gamma, \unit \seqar \sigma}{\vlhy{\Gamma \seqar \sigma}}
}
{
\vlhy{\Gamma , \N \seqar \sigma }
}
}
}
\]
    
    As before, it is routine to show that these reductions are derivable using $\cutrednorec$.
    
Now, specialising our simulation result to recursion on $\N$, we have  the  following regular coderivation for the recursor of system $\T$ (at type $\sigma$):
    \[
    % \begin{array}{rcl}
    %        \vliinf{\rec \sigma }{}{\Gamma, \N \seqar \sigma}{\Gamma \seqar \sigma}{\Gamma, \N, \sigma \seqar \sigma}  & \dfn & 
           \vlderivation{
            \vliin{\Ncnd }{}{\Gamma, \N \seqar \sigma}
            {
            \vlhy{\Gamma \seqar \sigma}
            }{
            \vliin{\cut}{}{\Gamma, \N \seqar \sigma}
              {
              \vlin{\Ncnd}{\bullet}{\Gamma, \N \seqar \sigma}{\vlhy{\vdots}}
              }{
              \vlhy{\Gamma, \N,\sigma  \seqar \sigma}
              }
            } 
           }
    % \end{array}
    \]
Indeed it is immediate that $\cmuLJ$ contains circular versions of system $\T$ from \cite{Das21:CT-preprint,Das21:CT-fscd,KuperbergPP21systemT}.
% , with $\Nzero$ beting $0$, $\Nsucc$ being $\succ$, and $\Ncnd$ being $\cnd$.
% \todo{better justify threading too}
\end{example}

\subsection{Reduction to the negative fragment}
\label{subsec:reduction-to-neg-frag}
It is folklore that coinductive types can be eliminated using inductive types (possibly at the loss of strict positivity) using, say, a version of the G\"odel-Gentzen negative translation, without affecting the class of representable functions (as long as $\N$ is included as a primitive data type) (see, e.g.,~\cite{avigad1998godel}).
Indeed this translation can be designed to eliminate other `positive' connectives too, in particular $+$.\footnote{Note that the attribution of `positive' or `negative' to a connective is unrelated to that of positive or negative context.}

The same trick does not quite work for coderivations since it introduces cuts globally that may break the progressing criterion in the limit of the translation.
\todonew{could elaborate here}
However a version of the Kolmogorov translation, more well behaved at the level of cut-free proof theory, is well suited for this purpose.
% 
In this section we establish such a reduction from $\cmuLJ$ to its `negative' fragment.
Not only is this of self-contained interest, being more subtle than the analogous argument for $\muLJ$ (and type systems with (co)inductive data types), this will also greatly simplify reasoning about the representable functions of $\cmuLJ$ in what follows, in particular requiring fewer cases in arguments therein.



\begin{defn}[Negative fragments] \label{defn:negative-fragments}
We define $\muLJneg$ as the subsystem of $\muLJ$ using only rules and cut-reductions over $\N, \times, \arrow, \mu$.
In particular we insist on the native rules and cut-reductions for $\N$ from \Cref{fig:native-n-rules+cutreds-muljnorec} to avoid extraneous occurrences of $+$ from $\N$ and remain internal to the fragment.
We define $\muLJnegnorec$ and $\cmuLJneg$ similarly, only as subsystems of $\muLJnorec$-coderivations and their cut-reductions.
\end{defn}


% \begin{remark}
% [Controlling sums]
% \todo{talk about inclusion of native rules for $\N$ to avoid occurrences of $+$ so that `using only $\N,\times,\arrow,\mu$' is well-defined.}
% \todo{somewhere need to say more about the $+$ and its subtlety for totality. Could foreshadow this already here.}
% \anupam{nb: in presence of proof normalisation the $+$ rule issue ceases to be a problem.}
% \end{remark}

The main result of this subsection is:
\begin{proposition}\label{prop:cmulj-into-cmuljminus}
    Any function on natural numbers representable in $\cmuLJ$ is also representable in $\cmuLJneg$.
\end{proposition}
\begin{proof}
[Proof idea]
We give a bespoke combination of a Kolmogorov negative translation and a Friedman-Dragalin `$A$-translation' (setting $A= \N$). We define the translations $\negtrans \cdot$ and $\unnegtrans \cdot$ from arbitrary types to types over $\{\N,\times, \to , \mu\}$ as follows, where $\lnot \sigma \dfn \sigma \arrow \N$:
\[
\begin{array}{r@{\ := \ }l}
\negtrans \sigma & \neg \unnegtrans \sigma\\ 
\unnegtrans X & \negn X \\
    \unnegtrans 1 & \N \\
    \unnegtrans{(\sigma \times \tau)} & \negn (\negtrans \sigma \times \negtrans \tau) \\
    \unnegtrans{(\sigma \to \tau)} & \neg( \negtrans\sigma \to \negtrans \tau) \\
    \unnegtrans{(\sigma+\tau)} & \neg   \negtrans \sigma \times \neg \negtrans \tau \\
    \unnegtrans{(\nu X \sigma)} & \neg \neg \mu X \negn \negtrans\sigma[\negn X / X] \\
    \unnegtrans{(\mu X \sigma)} &  \neg \mu X \negtrans{\sigma}
\end{array}
\]
The translation can be extended to coderivations by mapping every inference rule $\rrule$ to a gadget $\negtrans \rrule$  preserving threads.
Further details are given in~\Cref{subsec:proofs-of-cmulj-into-cmuljminus}.
\end{proof}


% \subsection{Further examples}
% \anupam{could revisit an example of stream computation in $\muLJ$, in particular tracking it through the simulation by $\cmuLJ$ and also the reduction to $\cmuLJneg$.}



\begin{example}\label{exmp:examples-of-coderivations}
The left coderivation of Figure~\ref{fig:examples-coderivations} shows the encoding of a  stream $n_0 :: n_1 :: n_2 \ldots$ by a (not necessarily regular) coderivation. 
Note that this coderivation is regular just if the stream is ultimately periodic.
The right coderivation shows the circular presentation of the  concatenation of a list and a stream into a stream  discussed  in Example~\ref{exmp:examples-of-derivations}. 
Note that, compared to the inductive encoding of this function, the circular one has an arguably more `explicit' computational meaning. 
Both coderivations are progressing, by the  \red{red} progressing threads in their  only infinite branches.

It is worth discussing  how   computation  over streams is simulated in $\cmuLJneg$ via the double negation translation illustrated in Proposition~\ref{prop:cmulj-into-cmuljminus}. 
% Recalling that $\neg \sigma \dfn \sigma  \to \N$,  
The type $S$ of streams is translated into $\neg \neg \neg \mu X \neg \neg \neg  (\negtrans \N \times \neg \neg \neg X)$, for some appropriate translation $\negtrans \N$ of the type for natural numbers. Hence,  computation over streams is simulated by computation over a type of the form  $( \sigma\to\N)\to \N$.
Note that this resembles (and embeds) the type $\N\arrow \N$ for representing streams in system $\T$, so in some sense we can see $\negtrans\cdot$-translation as extending/adapting the embedding of $S$ into $\N \arrow \N$.
\end{example}



\begin{figure}
    \centering
    \[
    \small
    \vlderivation{
\vlin{\rr \nu}{}{\seqar \red{S}}
{
\vliin{\rr \times}{}{\seqar \red{N \times S}}
{
\vltr{\cod {n_0}}{\seqar N}{\vlhy{\ }}{\vlhy{\ }}{\vlhy{\ }}
}
{
\vlin{\rr \nu}{}{\seqar \red{S}}
{
\vliin{\rr \times}{}{\seqar \red{N \times S}}
{
\vltr{\cod {n_1}}{\seqar N}{\vlhy{\ }}{\vlhy{\ }}{\vlhy{\ }}
}
{
\vlin{\rr \nu}{}{\seqar \red{S}}
{
\vlhy{\vdots}
}
}
}
}
}
}
\qquad
\vlderivation{
\vlin{\rr \to}{}{\red{L} \seqar S \to S}
{
\vlin{\lr \mu}{\bullet}{\red{L}, S \seqar S}
{
\vliin{\lr +}{}{\red{\unit +(N \times L}), S \seqar S}
{
\vlin{\lr \unit}{}{\unit, S \seqar S}{\vlin{\id}{}{S \seqar S}{\vlhy{}}}
}
{
\vlin{\rr \nu}{}{\red{N \times L}, S \seqar S}
{
\vlin{\lr \times}{}{\red{N \times L}, S \seqar N \times S}
{
\vliin{\rr \times}{}{N, \red{L}, S \seqar N \times S}
{
\vlin{\id}{}{N \seqar N}{\vlhy{}}
}
{
\vlin{\lr \mu  }{\bullet}{\red{L}, S \seqar S}
{
\vlhy{\vdots}
}
}
}
}
}
}
}
}
    \]
    \caption{Left: pointwise computation of a stream in $\muLJnorec$. Right: concatenation of a list and a stream in $\cmuLJ$.}
    \label{fig:list+stream-to-stream-cmuLJ}
     \label{fig:examples-coderivations}
\end{figure}